\documentclass[11pt]{article}
\usepackage[margin=2.3cm]{geometry}
\usepackage{amsmath,amssymb}
\usepackage{xr}
\externaldocument{main}
\usepackage[colorlinks=true,linkcolor=blue,urlcolor=blue]{hyperref}
\usepackage{enumitem}
\usepackage{microtype}
\setlength{\parskip}{0.45em}
\setlength{\parindent}{0pt}
\newenvironment{comment}[1]{\par\medskip\noindent\textbf{Comment~#1.}\ \itshape}{\par}
\newcommand{\resp}{\par\noindent\textbf{Response.}\ }

\title{Response to the reviewer (fourth round)\\[0.3em]\large A reliability-constrained, multi-objective
framework for minimum-mass design of steel frame structures under material and loading uncertainty}
\author{[Author name]}
\date{}

\begin{document}
\maketitle

We thank the reviewer for the recommendation of minor revision and for the precise list of remaining
corrections. No new computations were needed. All numerical results are unchanged from the third
revision. The changes concern wording, scope statements and two table entries. Section and table
numbers refer to the revised manuscript.

\begin{comment}{1 (status of the three-storey full-model estimate)}
Preserve the distinction between the sampled-strata result and the residual-adjusted result, and avoid
wording that implies a distribution-free bound for the complete full model.
\end{comment}
\resp We checked every place where the three-storey margin design is discussed (Abstract,
Section~\ref{sec:fullmc}, Table~\ref{tab:tail}, Section~\ref{sec:strict} and Conclusion~5).
\begin{itemize}[itemsep=0.1em]
\item Intervals and lower bounds that include the fitted $t_3$ residual are now called
\emph{model-conditional} confidence intervals and lower bounds wherever they appear. The same applies
to the caption of Table~\ref{tab:tail}.
\item Section~\ref{sec:strict} previously said that the target is ``met with 95\% confidence under the
$t_3$ residual model''. It now says that the model-conditional lower bound (3.43) exceeds the target
only under the assumed $t_3$ model, not under every tail model ($t_2$ gives 3.16), and that it is not
a distribution-free bound.
\item Conclusion~5 now reports the $t_2$ value next to the $t_3$ lower bound.
\end{itemize}

\begin{comment}{2 (code editions that define ``code-compliant'')}
Either update to the newer provisions or describe the study explicitly as a benchmark using AISC 360-16
and ASCE 7-05.
\end{comment}
\resp We retained the historical editions and now describe them explicitly as a fixed benchmark. We
did not update to the newer editions, because the wind-load statistics of the probabilistic model are
tied to the definition of the nominal wind load. From ASCE~7-10 onwards, wind loads are based on
strength-level wind speeds with a load factor of 1.0, so a change of edition would change the
nominal wind load and its statistics together, not just a load factor. The revision adds:
\begin{itemize}[itemsep=0.1em]
\item a ``Code basis'' paragraph in Section~\ref{sec:resistance}. It states that AISC~360-16 and ASCE~7-05
are used as a fixed, internally consistent benchmark, that both have been superseded (AISC~360-22 and
ASCE/SEI~7-22, now cited), and that the designs have not been checked against the current editions.
It also defines ``code-compliant'' as compliance with the AISC~360-16 checks under the ASCE~7-05
combinations;
\item the editions next to ``code-compliant'' in the Abstract, the list of formulations, the
Discussion and the preamble of the Conclusions;
\item in the code, a constant \texttt{CODE\_BASIS} in \texttt{rbo/limit\_states.py} that records the
editions implemented.
\end{itemize}

\begin{comment}{3 (number of validated designs)}
Table 1 still describes full second-order validation of eight designs; make the count consistent.
\end{comment}
\resp The ``This study'' row of Table~\ref{tab:literature} now reads ``full P-$\Delta$ MC of 10 designs
(8 principal + 2 three-storey variants)''. We also found that the caption of
Fig.~\ref{fig:validation} said ``eight designs'' although the figure shows ten. It now says ten
designs and names the two three-storey variants, and the note in Section~\ref{sec:fullmc} refers to
this figure as well.

\begin{comment}{4 (model factors in Table 3)}
Table 3 labels the model factors ``sensitivity only'', although they are also used in RBDO-C-MU.
\end{comment}
\resp The model-factor row of Table~\ref{tab:rv} is generated by \texttt{make\_figures.py}. It now
reads ``baseline: 0; sensitivity study: 0.05, 0.10; RBDO-C-MU (2 frames): 0.05''. This states the
c.o.v.\ for each use and agrees with Section~\ref{sec:probmodel} and Table~\ref{tab:mudesign}.

\begin{comment}{5 (minimum-mass claims and the genetic algorithm)}
Where global optimality has not been established, use wording such as best-found design.
\end{comment}
\resp
\begin{itemize}[itemsep=0.1em]
\item Section~\ref{sec:optimisation} now defines ``optimum'', ``optimal design'' and ``minimum-mass
design'' as the best feasible design found within the tested optimisation budget, not a proven global
minimum. It notes that ``minimum mass'' in the title names the objective of the problem.
\item In the Abstract, ``a least-mass design at a stated reliability'' became ``the lightest design
found, within the tested search budget, that meets a stated reliability''.
\item In the Discussion, ``generally the least-mass way to meet both targets'' became ``among the
designs compared, it generally gave the lowest mass that met both targets, with all masses being
best-found designs within the tested optimisation budget, not proven global minima''.
\item Two further statements were reworded in the same way: Section~\ref{sec:targets} (``the best
code-based design found is also the best reliability-based design found'') and Conclusion~6.
\end{itemize}

\begin{comment}{6 (mass baseline in the Abstract)}
\end{comment}
\resp The Abstract now reads ``2.7\% more steel than the point-estimate design (5.2\% more than the
deterministic optimum)''. Conclusion~5 states both premiums ($+2.7\%$ relative to RBDO-C and
$+5.2\%$ relative to DDO).

\section*{Minor editorial items}
\begin{itemize}[itemsep=0.2em]
\item \emph{Author information.} The placeholders remain in this version, because the author names
and affiliations have not yet been supplied to us. They will be inserted before formal submission.
\item \emph{Cross-references.} The manuscript was recompiled after the changes to
Tables~\ref{tab:literature} and~\ref{tab:rv}. There are no undefined references or citations.
\item \emph{Tail-model terminology.} ``Model-conditional confidence interval'' is now used throughout
(Comment~1).
\item \emph{Reproducibility.} The manuscript names the exact commit that produced it. The results are
unchanged from the third revision. A release tag and DOI will be created on acceptance.
\item \emph{Cost interpretation.} Section~\ref{sec:cost} and Conclusion~6 now state that the cost
index is an uncalibrated sensitivity model and has not been validated against actual project costs.
\item \emph{Confidence wording.} In the Abstract and Conclusion~3, ``met both targets with 95\%
confidence'' for the point-estimate designs now also names the baseline surrogate model.
\end{itemize}

\end{document}
